\documentclass[11pt,a4paper]{article}
\usepackage[margin=21mm,headheight=15pt]{geometry}
\usepackage[T1]{fontenc}
\usepackage{lmodern,amsmath,amssymb,booktabs,tabularx,enumitem,xcolor,fancyhdr}
\usepackage[hidelinks]{hyperref}
\definecolor{accent}{RGB}{0,114,178}
\setlength{\parindent}{0pt}
\setlength{\parskip}{6pt}
\setlist{nosep,leftmargin=*}
\pagestyle{fancy}
\fancyhf{}
\fancyhead[L]{\small Econometrics 25117 | Instructor guide}
\fancyhead[R]{\small 2026--2027}
\fancyfoot[L]{\small Private teaching notes}
\fancyfoot[R]{\thepage}
\newcommand{\heading}[1]{\section*{\color{accent}#1}}
\newcommand{\slideheading}[1]{\subsection*{\color{accent}#1}}
\newcommand{\cue}[2]{\par\noindent\textbf{#1}\quad #2\par}
\newcommand{\slideguide}[3]{\slideheading{Slide #1 | #2}\cue{Guide}{#3}}
\setlength{\emergencystretch}{1em}
\newcommand{\E}{\mathbb{E}}
\newcommand{\Var}{\operatorname{Var}}
\newcommand{\Cov}{\operatorname{Cov}}
\begin{document}
\begin{center}{\Large\bfseries Pooled data and first differences}\end{center}
\textbf{Slide route:} Lecture10v2.tex: Starting with Pooled Data through Limits of First-Differences. Stop before Fixed Effects Models.
\heading{Session 13 | Same units or new samples?}
\textbf{Outcomes:} Distinguish pooled cross-sections from panels, interpret time interactions, and show how first differences remove fixed unit characteristics.
\heading{100-minute route}
\begin{tabularx}{\textwidth}{@{}rX@{}}\toprule
0--10 & Retrieve interaction interpretation.\\
10--30 & Cross-section, pooled cross-section, fertility/gender examples.\\
30--45 & Time interactions and changing associations.\\
45--55 & Pause and identify the sampling unit.\\
55--75 & Panel structure and traffic-fatality examples.\\
75--90 & First differences for two and more periods.\\
90--100 & Limits and exit exercise.\\\bottomrule
\end{tabularx}
\heading{Slide-by-slide script}
\slideguide{1}{Title}{Open by asking whether observations from two dates are the same people. That determines the data structure.}
\slideguide{2}{Roadmap}{Preview cross-sections, pooled data, interactions over time, panels, and first differences.}
\slideguide{3}{Last lesson}{Retrieve the IV distinction between a new source of variation and a new estimator.}
\slideguide{4}{Starting with pooled data}{Name the unit, time index, outcome, and sample before interpreting any pooled regression.}
\slideguide{5}{Cross-section}{Explain that one cross-section observes many units at one time.}
\slideguide{6}{Pooled cross-section}{Contrast repeated independent samples with repeated observations of the same people.}
\slideguide{7}{Pooled cross-section}{Use a time dummy to absorb a common average shift, but warn that it does not create a panel.}
\slideguide{8}{Pooled cross-section}{Ask whether the composition of the sample can change over time and why that matters.}
\slideguide{9}{Fertility over time}{Identify the outcome, cohort or time variable, and the reference category.}
\slideguide{10}{Fertility over time}{Interpret a time indicator as a conditional difference, not automatically an individual change.}
\slideguide{11}{Fertility over time}{Ask what omitted cohort characteristics could move with time.}
\slideguide{12}{Fertility over time}{Use the interaction to let a relationship differ across periods.}
\slideguide{13}{Fertility over time}{Calculate one fitted value from the displayed specification and identify the baseline group.}
\slideguide{14}{Fertility over time}{Distinguish a time trend from a causal treatment effect.}
\slideguide{15}{Gender gap}{Write the group contrast at the reference date and explain its units.}
\slideguide{16}{Gender gap}{Use the interaction to calculate how the gap changes later.}
\slideguide{17}{Gender gap}{Ask whether the later gap difference is causal without a parallel-trends argument.}
\slideguide{18}{Testing differences}{Frame the changing-gap question as a joint or linear restriction when appropriate.}
\slideguide{19}{Transition to panel data}{Explain why following the same units changes what can be removed.}
\slideguide{20}{Panel data}{Define a panel using unit and time indices. Have students label both in the traffic-fatality example.}
\slideguide{21}{Panel data}{Contrast between-unit variation with within-unit changes.}
\slideguide{22}{Why panel data?}{Explain that fixed characteristics can be differenced out, while time-varying confounders remain.}
\slideguide{23}{Traffic fatalities, 1982}{Read one state's outcome and beer tax, then identify the comparison being made.}
\slideguide{24}{Traffic fatalities, 1983}{Ask students to compute the change for one state.}
\slideguide{25}{Traffic fatalities, 1984}{Continue the table and distinguish levels from changes.}
\slideguide{26}{Traffic fatalities, 1985}{Ask whether a common time shock would disappear from first differences.}
\slideguide{27}{Traffic fatalities, 1986}{Keep state labels fixed; this is a panel, not a new cross-section.}
\slideguide{28}{Traffic fatalities, 1987}{Use one state as a numerical check on the differencing arithmetic.}
\slideguide{29}{Traffic fatalities, 1988}{Ask what information is lost when unit means are removed.}
\slideguide{30}{First differences}{Write $\Delta Y_{it}=\beta\Delta X_{it}+\Delta u_{it}$ and show the fixed effect cancels.}
\slideguide{31}{First differences}{Use the three-unit example from this guide to calculate the slope through the origin with a common change.}
\slideguide{32}{First differences when $T>2$}{Explain that successive differences share level errors and can be serially correlated.}
\slideguide{33}{Limits of first differences}{List time-varying confounding, measurement error amplification, and the need for repeated observations.}
\slideguide{34}{Limits of first differences}{Ask whether independent repeated cross-sections can be differenced person by person. No.}
\slideguide{35}{Transition}{Close by previewing within transformations and fixed effects for the next class.}

\heading{Side example | Pooled observations}
\cue{Use with slides}{Lecture10v2 slides 5--18: Cross-Section, Pooled Cross-Section, Fertility over time, Gender Gap, and Testing Differences in Pooled regressions. Use the two-by-two interaction example while moving through slides 16--18.}
In a synthetic pooled wage model,
\[
wage=10+2D_{\rm later}+3D_{\rm group}
-1(D_{\rm later}D_{\rm group})+u,
\]
the group gap is 3 initially and 2 later. The interaction is the change in the gap, $-1$. With new people sampled at each date, this is not the average change in the same individuals' wages.
\cue{Ask}{Does a two-by-two interaction automatically identify a causal difference-in-differences effect? No. That reading requires a design and additional assumptions such as a suitable parallel-trends counterfactual. Here the example is an interpretation exercise.}
\cue{Panel definition}{A panel repeatedly observes the same units. Time labels on independent cross-sections do not turn them into a panel. Have students name the unit, time index, and outcome before interpreting any coefficient.}

\newpage
\heading{Worked example | Removing a fixed effect}
\cue{Use with slides}{Lecture10v2 slides 20--34: panel data through Limits of First-Differences. Use the three-unit table when the deck reaches Traffic Fatalities first differences, especially slides 31--34.}
Suppose
\[
Y_{it}=\alpha_i+\beta X_{it}+\lambda_t+u_{it}.
\]
With two periods,
\[
\Delta Y_i=\beta\Delta X_i+\Delta\lambda+\Delta u_i.
\]
The unit effect $\alpha_i$ cancels. A common time shock remains and can be captured by an intercept in the two-period differenced regression.
\begin{center}\begin{tabular}{lrrrr}\toprule
Unit & $X_1$ & $X_2$ & $Y_1$ & $Y_2$\\\midrule
A & 1 & 2 & 10 & 13\\
B & 2 & 4 & 20 & 25\\
C & 3 & 6 & 30 & 37\\\bottomrule
\end{tabular}\end{center}
Differences are $(\Delta X,\Delta Y)=(1,3),(2,5),(3,7)$. The fitted differenced relationship is
\[
\Delta\widehat Y=1+2\Delta X.
\]
\cue{Ask}{What would happen if we forced this line through the origin? The slope would be $(3+10+21)/(1+4+9)=34/14\simeq2.429$. Omitting the common change contaminates this synthetic comparison.}
\cue{Identification comment}{Removing fixed attributes does not remove time-varying confounders. If a district changes tax policy in response to an unobserved changing risk, first differences need not identify the causal effect.}
\cue{More than two periods}{Successive differences share errors: $\Delta u_t=u_t-u_{t-1}$ and $\Delta u_{t-1}=u_{t-1}-u_{t-2}$. Even independent level errors induce dependence between neighboring differenced errors.}
\cue{Exit answers}{Time-invariant regressors disappear under differencing. A common time shock does not disappear unless it is unchanged. Repeated independent samples cannot generally be differenced person by person.}
\cue{If short / preparation}{Show selected years of the traffic-fatality figures rather than each in detail. Preserve the three-unit example. Next time extend removal of fixed attributes to all periods using within transformations.}
\textbf{Source:} Lecture10v2.tex, pooled-data and first-difference sections; synthetic numerical table.

\end{document}
